SLSC は, 確率紙上の点列と理論線の離れ方を一つの数値で表す指標であり, 視覚的な判断と対応しやすい.
本稿の結果は, この診断量としての有用性を踏まえた上で, 分布選択の数値基準として用いる際の留意点を整理するものである.

\rev{第一に, 固定閾値 \(0.04\) は従来の実務上の目安として用いられてきたが, 真の分布を当てはめた場合でも SLSC は標本サイズとともに低下するため, 標本サイズや候補分布群に依存しない一律の判定境界として用いることは適切ではない.}
\rev{この点は, 林ら\cite{HayashiEtAl2012}が SLSC を検定統計量として扱い, 標本サイズや想定分布に応じた棄却限界値を用いる考え方とも整合する.}

第二に, 異なる分布を比較する際には, SLSC をどの空間で計算しているかが本質的に重要である.
線形標準化であれば標準変量への変換は SLSC の値を変えないが, LP3 や LN3 のような非線形標準化では, 変換後の空間で測った距離は元の水文量空間で測った距離とは一致しない.
LN3(S) が小さい SLSC を与えやすいという結果は, LN3 そのものの有用性を否定するものではなく, 非線形標準化を介した \(S\) 空間の SLSC を他分布の SLSC と同じ物差しで比較できるかという問題を示している.
分布間の適合度を共通尺度として比較したい場合には, \(X\) 空間で評価するか, 標準化の違いが結果に与える影響を併記することが望ましい.
\rev{したがって, 分布間比較を主目的とする場合には, 非線形標準化の影響を受ける \(S\) 空間 SLSC ではなく, 共通尺度としての \(X\) 空間 SLSC を主指標として用いることの検討が重要である.}

\rev{第三に, SLSC+ジャックナイフ法では, 最終的にはジャックナイフ幅 \(W_J\) の最小値で分布を選ぶものの, その前段で SLSC によるスクリーニングを行うため, SLSC の候補分布依存性が採択結果に残る.
したがって, LN3(S) が小さい SLSC を与えやすい場合には, LN3 が候補集合に残りやすく, その影響は \(W_J\) による最終選択にも及びうる.
一方, AIC ではこのような SLSC スクリーニングを経ないため, 少なくとも本稿の計算条件では LN3 への一方向的な採択集中はみられなかった.
ただし, AIC も \(N<150\) の範囲で常に真の分布を選択するわけではなく, GEV, LP3, LN3 では SqrtEt への採択が多いなど, 採択結果には別の偏りが残った.
このため, AIC は SLSC 系規準の特徴を相対化する参照指標として有用であるが, それ自体を無条件に正しい分布選択基準とみなすことはできない.
なお, GEV, LP3, Exp, LN3 などの端点が未知母数に依存する分布では最尤推定の正則条件が自明ではないため, AIC の解釈にはこの点でも注意を要する.}

\rev{さらに, ジャックナイフ幅 \(W_J\) は, 候補分布および推定手順を所与としたときに, 対象とする確率水文量の推定値が標本の1点削除に対してどの程度変動するかを表す量である. 
すなわち, \(W_J\) は推定安定性の指標であり, 分布への当てはまりやモデル誤指定に起因するバイアスの小ささを直接評価するものではない. 
したがって, \(W_J\) の小ささのみを根拠として当該分布を真の母分布に近いと判断することはできない.}

以上を踏まえると, SLSC は, 確率紙と対応した直感的な診断量として用いることが適切である.
分布選択の根拠として用いる場合には, 標本サイズ, 標準変量の定義, \(X\) 空間と \(S\) 空間の違い, および候補分布群への依存性を明示する必要がある.
このような整理を行うことで, 従来の実務で蓄積されてきた SLSC の知見を活かしつつ, 分布比較に伴う解釈上の混乱を小さくできると考えられる.

\rev{本稿にはいくつかの限界もある.
真の分布の母数は京都地点の降雨データを参考に設定しており, 歪度係数等の統計特性が異なる地点でも同様の傾向が成り立つかは今後確認する必要がある.
また, ジャックナイフ幅は50年再現期間の確率水文量に基づいており, 他の再現期間では採択傾向が変わりうる.}
3母数分布の尤度や推定安定性を含めた分布選択法の整理は, 今後の課題である.
